\providecommand{\weakconv}{\rightharpoonup}
\providecommand{\K}{\mathbb{K}}

\subsection{Слабая сходимость и линейные ограниченные операторы. Слабо ограниченные множества. Теорема Хана}

\begin{theorem}[Слабая сходимость и ограниченный линейный оператор]
	Пусть \(E_1,E_2\) --- нормированные пространства,
	\[
	A\in \mathcal{L}(E_1,E_2).
	\]
	Если
	\[
	x_n\weakconv x
	\qquad \text{в } E_1,
	\]
	то
	\[
	Ax_n\weakconv Ax
	\qquad \text{в } E_2.
	\]
\end{theorem}

\begin{proof}
	Докажем утверждение прямо по определению слабой сходимости.
	
	\textbf{1. Сведение к функционалу на \(E_1\).}
	
	Возьмём произвольный функционал
	\[
	g\in E_2^*.
	\]
	Чтобы доказать слабую сходимость
	\[
	Ax_n\weakconv Ax
	\quad \text{в } E_2,
	\]
	нужно проверить, что
	\[
	g(Ax_n)\to g(Ax).
	\]
	
	Рассмотрим композицию
	\[
	f=g\circ A,
	\qquad
	f(y)=g(Ay),
	\qquad
	y\in E_1.
	\]
	Тогда \(f\) линеен как композиция линейных отображений. Кроме того, для любого \(y\in E_1\)
	\[
	|f(y)|
	=
	|g(Ay)|
	\leqslant
	\|g\|\cdot\|Ay\|
	\leqslant
	\|g\|\cdot\|A\|\cdot\|y\|.
	\]
	Следовательно,
	\[
	f\in E_1^*.
	\]
	
	\textbf{2. Использование слабой сходимости в \(E_1\).}
	
	Так как
	\[
	x_n\weakconv x
	\qquad \text{в } E_1,
	\]
	то для функционала \(f=g\circ A\in E_1^*\) имеем
	\[
	f(x_n)\to f(x).
	\]
	То есть
	\[
	(g\circ A)(x_n)\to (g\circ A)(x),
	\]
	или, что то же самое,
	\[
	g(Ax_n)\to g(Ax).
	\]
	
	Функционал \(g\in E_2^*\) был произвольным, значит
	\[
	\forall g\in E_2^*
	\qquad
	g(Ax_n)\to g(Ax).
	\]
	По определению слабой сходимости это означает
	\[
	Ax_n\weakconv Ax.
	\]
\end{proof}

\begin{definition}[Слабо ограниченное множество]
	Пусть \(E\) --- нормированное пространство. Множество
	\[
	S\subset E
	\]
	называется \textbf{слабо ограниченным}, если для любого функционала
	\(f\in E^*\) множество
	\[
	f(S)=\{f(x):x\in S\}
	\]
	ограничено в \(\K\).
	
	Иными словами,
	\[
	\forall f\in E^*
	\quad
	\exists C(f)>0
	\quad
	\forall x\in S
	\qquad
	|f(x)|\leqslant C(f).
	\]
	Важно: константа \(C(f)\) может зависеть от функционала \(f\), но не зависит от точки \(x\in S\).
\end{definition}

\begin{remark}[Сравнение с обычной ограниченностью]
	Обычная ограниченность множества \(S\subset E\) означает:
	\[
	\exists C>0
	\quad
	\forall x\in S
	\qquad
	\|x\|\leqslant C.
	\]
	
	Слабая ограниченность означает:
	\[
	\forall f\in E^*
	\quad
	\exists C(f)>0
	\quad
	\forall x\in S
	\qquad
	|f(x)|\leqslant C(f).
	\]
	
	Теорема Хана утверждает, что эти два условия на самом деле эквивалентны.
\end{remark}

\begin{theorem}[Хана]
	Пусть \(E\) --- нормированное пространство и \(S\subset E\). Тогда
	\[
	S \text{ слабо ограничено}
	\quad\Longleftrightarrow\quad
	S \text{ ограничено}.
	\]
\end{theorem}

\begin{proof}
	Докажем обе импликации.
	
	\textbf{1. Ограниченность \(\Rightarrow\) слабая ограниченность.}
	
	Пусть \(S\) ограничено. Тогда существует \(R>0\), такое что
	\[
	\|x\|\leqslant R
	\qquad
	\forall x\in S.
	\]
	Возьмём произвольный функционал \(f\in E^*\). Тогда для любого \(x\in S\)
	\[
	|f(x)|
	\leqslant
	\|f\|\cdot\|x\|
	\leqslant
	\|f\|\cdot R.
	\]
	Значит, множество \(f(S)\) ограничено. Так как \(f\in E^*\) был произвольным,
	то \(S\) слабо ограничено.
	
	\textbf{2. Слабая ограниченность \(\Rightarrow\) ограниченность.}
	
	Докажем от противного. Пусть \(S\) слабо ограничено, но не ограничено.
	
	\begin{itemize}
		\item \textbf{Выберем быстро растущую последовательность.}
		
		Из неограниченности \(S\) следует, что для каждого \(n\in\mathbb{N}\)
		можно выбрать
		\[
		x_n\in S
		\]
		так, что
		\[
		\|x_n\|>n^2.
		\]
		Положим
		\[
		y_n=\frac{x_n}{n}.
		\]
		
		\item \textbf{Покажем, что \(y_n\weakconv 0\).}
		
		Возьмём произвольный функционал \(f\in E^*\). Так как \(S\) слабо ограничено,
		то существует \(C(f)>0\), такое что
		\[
		|f(x)|\leqslant C(f)
		\qquad
		\forall x\in S.
		\]
		В частности,
		\[
		|f(x_n)|\leqslant C(f).
		\]
		Тогда
		\[
		|f(y_n)|
		=
		\left|f\left(\frac{x_n}{n}\right)\right|
		=
		\frac{1}{n}|f(x_n)|
		\leqslant
		\frac{C(f)}{n}
		\to 0.
		\]
		Следовательно,
		\[
		f(y_n)\to 0=f(0)
		\qquad
		\forall f\in E^*.
		\]
		По определению слабой сходимости
		\[
		y_n\weakconv 0.
		\]
		
		\item \textbf{Получим противоречие с ограниченностью слабо сходящейся последовательности.}
		
		По лемме из предыдущего билета всякая слабо сходящаяся последовательность ограничена.
		Значит, последовательность \(\{y_n\}\) должна быть ограничена.
		
		Но по построению
		\[
		\|y_n\|
		=
		\left\|\frac{x_n}{n}\right\|
		=
		\frac{1}{n}\|x_n\|
		>
		\frac{1}{n}\cdot n^2
		=
		n.
		\]
		Следовательно,
		\[
		\|y_n\|>n
		\qquad
		\forall n\in\mathbb{N}.
		\]
		То есть \(\{y_n\}\) не ограничена. Противоречие.
	\end{itemize}
	
	Значит, предположение о неограниченности \(S\) неверно. Следовательно,
	\[
	S \text{ ограничено}.
	\]
\end{proof}

